% Start of the text from the radiation safety manual in Spanish with spelling corrections where necessary.
Una característica desfavorable de los dosímetros de película, es que la dosis se puede determinar con un error de 20 a 30%.

\section*{Alarma sonora}

La alarma sonora tiene como característica el que anuncian en forma audible que se está en presencia de un campo de radiación alta, estos monitores están calibrados de tal forma que cuando están en presencia de un campo de radiación de 1mR/hr emiten un "bip" cada 30 segundos.

Así, si un trabajador está en un campo de radiación de 2mR/hr, su alarma dará 2 "bip"s cada 30 segundos.

En el caso de que un trabajador sin saber estuviese colocado en un campo de radiación alto, digamos 60mR/hr, su alarma estuviera emitiendo 60 "bip"s cada 30 segundos, es decir, 2 cada segundo.

Estos monitores de alarma tienen dos posiciones, en la baja (que significa baja sensibilidad) es donde se usan la mayoría parte del tiempo. En la alta sensibilidad comportan de similar, solo que emite 30 "bip"s en 30 segundos en un campo de 1 mR/hr, esta escala se usa solo para inspeccionar paquetes.

\section*{Comparación}

En resumen podemos mencionar las ventajas y desventajas de cada uno:

% Dosimetro de lectura directa
\textbf{DOSÍMETRO DE LECTURA DIRECTA.}
\begin{itemize}
    \item Ventajas: Información rápida y a la mano.
    \item Desventajas: Si se golpea puede perder la información. Si se recibe mucha dosis se satura y pierde la lectura.
\end{itemize}

% Dosimetro de pelicula
\textbf{DOSÍMETRO DE PELÍCULA.}
\begin{itemize}
    \item Ventajas: No se satura y no pierde la información al golpearse.
    \item Desventajas: Necesita revelarse para ser leído. Es dañado por la humedad.
\end{itemize}

% Monitor de alarma
\textbf{MONITOR DE ALARMA.}
\begin{itemize}
    \item Ventajas: Alerta en forma audible si el nivel de radiación es alto.
    \item Desventajas: No proporciona una lectura del nivel de radiación.
\end{itemize}
% End of manual text extraction and LaTeX conversion.
